%%%%%%%%%%%%%%%%%%%%%%%%%%%%%%%%%%%%%%%%%%%%%%%%%%%%%%%%%%%%%%%%%%%%%%%%%%%%%%%
% Academic CV — Olivier Bisson
% Single-file moderncv template (clean, academic-oriented)
%%%%%%%%%%%%%%%%%%%%%%%%%%%%%%%%%%%%%%%%%%%%%%%%%%%%%%%%%%%%%%%%%%%%%%%%%%%%%%%

\documentclass[11pt,a4paper,sans]{moderncv}

% Style
\moderncvstyle[nosymbols]{classic} % portable letter labels for contact details
\moderncvcolor{blue}         % options: blue, orange, green, red, purple, grey, black

% Encoding
\usepackage[utf8]{inputenc}
\usepackage[T1]{fontenc}
\usepackage{lmodern}

% Layout
\usepackage[scale=0.90]{geometry}
\setlength{\hintscolumnwidth}{2.7cm}

% Links
\AtBeginDocument{\hypersetup{hidelinks}}

%----------------------------------------------------------------------------------------
%  NAME AND CONTACT
%----------------------------------------------------------------------------------------

\name{Olivier}{Bisson}
\title{PhD Student in Geometric Statistics}

% Institutional address (recommended for academic CV)
\address{Inria (EPIONE), Universit\'e C\^ote d'Azur}{2004 Route des Lucioles, 06560 Valbonne, France}

% Optional: phone (comment out if you prefer)
% \phone[mobile]{+33~6~12~35~34~97}

\email{olivier.bisson@inria.fr}

% Optional: website / scholar (fill later)
\homepage{olivierbisson.github.io}
% \social[github]{yourhandle}
% \social[googlescholar]{YOURID}{Google Scholar}

% Optional photo
% \photo[70pt][0.4pt]{me}

%----------------------------------------------------------------------------------------

\begin{document}

\makecvtitle

%----------------------------------------------------------------------------------------
%  SHORT PROFILE
%----------------------------------------------------------------------------------------

\section{Profile}
\cvitem{}{I study the \textbf{geometry of smooth manifolds and stratified spaces}, with a focus on correlation and covariance matrices. My work combines \textbf{Riemannian and quotient geometry}, \textbf{orbit-type stratification}, and geometric statistics. I investigate intrinsic distances, rank-constrained geodesics, curvature, holonomy, and volume measures on spaces of low-rank correlation matrices, alongside log-Euclidean constructions for full-rank matrix manifolds. Applications include statistical modeling of brain connectivity trajectories.}

%----------------------------------------------------------------------------------------
%  RESEARCH INTERESTS
%----------------------------------------------------------------------------------------

\section{Research interests}
\cvitem{}{Geometry of stratified spaces and orbit-type quotients; Riemannian submersions and principal bundles; intrinsic metrics and rank-constrained geodesics; curvature, connections, and holonomy; volume measures on matrix spaces; log-Euclidean geometry; geometric statistics and applications to neuroimaging.}

%----------------------------------------------------------------------------------------
%  CURRENT POSITION
%----------------------------------------------------------------------------------------

\section{Current position}
\cventry{Oct. 2023 -- present}{PhD Student in Geometric Statistics}{Inria (EPIONE) \& Universit\'e C\^ote d'Azur}{France}{}{\textbf{Thesis:} \emph{Riemannian geometry and orbit type stratification of correlation matrices}.\\
\textbf{Advisor:} Xavier Pennec.\\
\textbf{PhD defense:} scheduled for 12 November 2026.\\
\textbf{Funding:} 3IA (Interdisciplinary Institute for Artificial Intelligence).}

%----------------------------------------------------------------------------------------
%  EDUCATION
%----------------------------------------------------------------------------------------

\section{Education}
\cventry{2023}{M.Sc. in Mathematical Engineering}{Stochastic and Statistical Modeling}{Universit\'e C\^ote d'Azur, France}{Highest honors}{}
\cventry{2022}{M.Sc. in Pure and Applied Mathematics}{Algebra and Geometry}{Universit\'e C\^ote d'Azur}{France}{}
\cventry{2020}{B.Sc. in Fundamental Mathematics}{}{Universit\'e C\^ote d'Azur}{France}{}

%----------------------------------------------------------------------------------------
%  PUBLICATIONS
%----------------------------------------------------------------------------------------

\section{Published journal papers}

\cvitem{2024}{\textbf{O. Bisson}, X. Pennec.
\emph{Differential of the Stretch Tensor for Any Dimension with Applications to Quotient Geodesics}.
\emph{Comptes Rendus. Math\'ematique}, 362 (2024), pp. 1847--1856.
\href{https://doi.org/10.5802/crmath.692}{doi:10.5802/crmath.692}.
\href{https://comptes-rendus.academie-sciences.fr/mathematique/item/10.5802/crmath.692.pdf}{PDF}.}

\section{Conference papers}

\cvitem{2025}{\textbf{O. Bisson}, Y. Aeschlimann, S. Deslauriers-Gauthier, X. Pennec.
\emph{Log-Euclidean Frameworks for Smooth Brain Connectivity Trajectories}.
In: \emph{Geometric Science of Information (GSI)}, St Malo, France, 2025, pp. 194--204.
Springer. \href{https://doi.org/10.1007/978-3-032-03924-8_20}{doi:10.1007/978-3-032-03924-8\_20}.
\href{https://arxiv.org/abs/2507.15374}{arXiv:2507.15374}.
\textit{Nominated for Best Paper Award.}}

\section{Preprints}

\cvitem{2026}{\textbf{O. Bisson}, D. Loiseaux, X. Pennec.
\emph{On the Geometry of Low Rank Correlation Matrices}.
\href{https://arxiv.org/abs/2610.28475}{arXiv:2610.28475}.}

\cvitem{2025}{\textbf{O. Bisson}, X. Pennec. \emph{Log-Euclidean Lie Groups}.
\href{https://arxiv.org/abs/2511.13380}{arXiv:2511.13380}.}

%----------------------------------------------------------------------------------------
%  TALKS
%----------------------------------------------------------------------------------------

\clearpage
\section{Talks \& conferences}
\cvitem{Oct. 2026}{\emph{On the Geometry of Low Rank Correlation Matrices}. \href{https://www.imsi.institute/activities/statistical-foundations-for-the-analysis-of-non-euclidean-data/}{Statistical Foundations for the Analysis of Non-Euclidean Data}, Institute for Mathematical and Statistical Innovation (IMSI), University of Chicago, Chicago, USA, 21 October 2026.}
\cvitem{Jul. 2026}{\emph{Poster: Log-Euclidean Lie Groups and rs-fMRI Brain Connectivity Trajectories}. Foundations of Computational Mathematics (FoCM), Faculty of Mathematics, University of Vienna, Vienna, Austria, 8--18 July 2026.}
\cvitem{May 2026}{\emph{Log-Euclidean Geometry of Full-Rank Correlation Matrices}. Geometry, Topology and Machine Learning 2026 (GML 2026), Sorbonne Universit\'e, Paris, France, 18--22 May 2026.}
\cvitem{Dec. 2025}{\emph{Mini-course: Quotient Spaces and Related Differential Geometry}. EPIONE Seminar, Inria Sophia Antipolis, France, 12 December 2025.}
\cvitem{Oct. 2025}{\emph{Log-Euclidean Frameworks for Smooth Brain Connectivity Trajectories}. Geometric Science of Information (GSI), St Malo, France, 29 October 2025. \textit{Nominated for Best Paper Award.}}
\cvitem{Mar. 2025}{\emph{Geometry, Stratification and Applications of Structured Correlation Matrices}. EPIONE Seminar, Inria Sophia Antipolis, France, 5 March 2025.}
\cvitem{Feb. 2025}{\emph{Differentiable Structures on Correlation Matrices and Applications in Neuroimaging}. \emph{Infinite-dimensional Geometry: Theory and Applications}, ESI (Vienna), Austria, 10 February 2025.}
\cvitem{Dec. 2024}{\emph{Geometry, Stratification and Applications of Structured Correlation Matrices}. 3IA C\^ote d'Azur PhD Seminar, Nice, France, 6 December 2024.}
\cvitem{Nov. 2024}{\emph{Riemannian Metrics on Correlation Matrices with Applications in Neuroimaging}. Mathematical Imaging and Random Geometry (RT MAIAGES), Nice, France, 25 November 2024.}
\cvitem{Oct. 2024}{\emph{Introduction to the Riemannian Geometry of Correlation Matrices}. EPIONE Seminar, Inria Sophia Antipolis, France, 18 October 2024.}
\cvitem{Feb. 2024}{\emph{A framework to put metrics on quotient spaces, with application in medical image analysis}. EPIONE Annual Retreat, Auron, France.}
\cvitem{Dec. 2023}{\emph{Poster: Geometry, stratification and applications of structured correlation matrices}. 3IA Scientific Day, SKEMA, Sophia Antipolis, France.}
\cvitem{Oct. 2023}{\emph{A framework to put metrics on quotient spaces, with application in medical image analysis}. EPIONE Seminar, Inria Sophia Antipolis, France.}

%----------------------------------------------------------------------------------------
%  AWARDS AND FUNDINGS
%----------------------------------------------------------------------------------------

\section{Awards \& Funding}
\cvitem{2023--2026}{\textbf{3IA Côte d’Azur PhD Funding}, Université Côte d’Azur / Inria.}
\cvitem{2025}{\textbf{Best Paper Award nomination} (GSI 2025) for \emph{Log-euclidean frameworks for smooth brain connectivity trajectories}.}


%----------------------------------------------------------------------------------------
%  TEACHING
%----------------------------------------------------------------------------------------

\section{Teaching}
\cventry{2025--2026}{Teaching assistant (TD)}{Introduction \`a l'analyse (L1)}{Universit\'e C\^ote d'Azur}{}{Sets and functions; limits and continuity; differentiability; function study; injective/surjective/bijective maps; integrals and primitives.}
\cventry{2024--2025}{Teaching assistant (TD)}{Compl\'ements d'alg\`ebre lin\'eaire (L2)}{Universit\'e C\^ote d'Azur}{}{Orthogonality and projections; Gram--Schmidt; least squares; isometries; introduction to PCA.}
\cventry{2024--2025}{Teaching assistant (TD)}{Calculus II (L2)}{Universit\'e C\^ote d'Azur}{}{Complex numbers; linear ODEs; multivariable calculus (partials, gradient, tangent planes, level sets, local extrema).}
\cventry{2023--2024}{Teaching assistant (TD)}{Maths Approfondissement 2 (L1)}{Universit\'e C\^ote d'Azur}{}{Exponential/logarithm refreshers; linear ODEs; Riemann integration basics; partial fractions; Euclidean vector geometry.}

%----------------------------------------------------------------------------------------
%  SOFTWARE / OPEN SOURCE
%----------------------------------------------------------------------------------------

\clearpage
\section{Software \& open-source}
\cvitem{}{Contributor to the \href{https://github.com/geomstats/geomstats}{\texttt{geomstats}} Python package, with implementations related to Riemannian geometry on SPD and correlation matrix manifolds.}

%----------------------------------------------------------------------------------------
%  PREVIOUS RESEARCH EXPERIENCE (SELECTED)
%----------------------------------------------------------------------------------------

\section{Previous Research experiences}
\cventry{Apr. 2023 -- Oct. 2023}{Research intern}{Inria (EPIONE), Universit\'e C\^ote d'Azur}{Nice/Valbonne, France}{}{\textbf{Topic:} \emph{Interpolation and extrapolation of correlation matrices for dynamic functional connectivity}.\\
\textbf{Supervisors:} Xavier Pennec, Samuel Deslauriers-Gauthier.\\
Implementation and evaluation of Riemannian metrics on correlation matrices for geodesic interpolation and extrapolation on resting-state fMRI correlation time-series.}

\cventry{Sep. 2022 -- Feb. 2023}{Research project}{LJAD \& CEPAM, Université Côte d'Azur}{Nice, France}{}{%
\textbf{Topic:} \emph{Supervised classification of wood charcoal fragments: from machine to active learning}.\\
\textbf{Supervisor:} Marco Corneli.\\
Development of learning methods to automatically identify wood charcoal remains from a Southern-African context (M.Sc. capstone project).}

\cventry{Mar. 2022 -- Jun. 2022}{Research intern}{LJAD, Université Côte d'Azur}{Nice, France}{}{%
\textbf{Topic:} \emph{Quantum representations of mapping class groups}.\\
\textbf{Supervisor:} Jeremy Toulisse.\\
Topological/combinatorial approach (skein theory) to construct quantum representations of mapping class groups.}


% Optional older items (uncomment if you want to keep them)
% \cventry{Sep. 2022 -- Feb. 2023}{Research project}{LJAD \& CEPAM, Universit\'e C\^ote d'Azur}{Nice, France}{}{\emph{Supervised classification of wood charcoal fragments: from machine to active learning}.}
% \cventry{Mar. 2022 -- Jun. 2022}{Research intern}{LJAD, Universit\'e C\^ote d'Azur}{Nice, France}{}{\emph{Quantum representations of mapping class groups}. Supervisor: Jeremy Toulisse.}

%----------------------------------------------------------------------------------------
%  LANGUAGES
%----------------------------------------------------------------------------------------

\section{Languages}
\cvitemwithcomment{French}{Native}{Fluent}
\cvitemwithcomment{English}{Fluent}{Reading, writing, speaking}
\cvitemwithcomment{Italian}{Basic}{Working knowledge}

%----------------------------------------------------------------------------------------
%  OPTIONAL: INTERESTS (SHORT)
%----------------------------------------------------------------------------------------

\section{Additional}
\cvitem{}{Outside research: Volleyball, rock climbing, sailing, beach volleyball and swimming. I am deeply convinced that \textbf{anyone} can do mathematics.}

\end{document}
